\unitchapter{2}{10}{Artificial Intelligence}{p2-10}

\unitmeta{Expected questions: 8--10 · Target: 8+ · Type: search traces,
alpha-beta, fuzzy operations, perceptron, concepts (Rich \& Knight,
Russell \& Norvig)}

\begin{sylbox}

\begin{itemize}
\tightlist
\item[$\square$]
  Approaches to AI: Turing test, rational agents, state-space representation, heuristic search, game playing, minimax, alpha-beta pruning
\item[$\square$]
  Knowledge representation: logic, semantic networks, frames, rules, scripts, conceptual dependency, ontologies, expert systems, uncertainty
\item[$\square$]
  Planning: components, linear \& non-linear planning, goal stack, hierarchical planning, STRIPS, partial-order planning
\item[$\square$]
  NLP: grammars, parsing techniques, semantic analysis, pragmatics
\item[$\square$]
  Multi-agent systems: agents vs objects vs expert systems, MAS structure, semantic web, agent communication, ontologies, tools
\item[$\square$]
  Fuzzy sets: membership functions, fuzzification/defuzzification, operations, linguistic variables, fuzzy relations, fuzzy inference \& control
\item[$\square$]
  Genetic algorithms: encoding, operators, fitness function, GA cycle
\item[$\square$]
  Neural networks: supervised, unsupervised, reinforcement learning; perceptron, MLP, SOM, Hopfield
\end{itemize}

\end{sylbox}

\needspace{130pt}

\hypertarget{p2-10--1-approaches-to-ai}{%
\section{Approaches to AI}\label{p2-10--1-approaches-to-ai}}

\begingroup\small

\par\medskip\noindent\hfil\begin{tabular}{@{}
  >{\raggedright\arraybackslash}p{(\columnwidth - 4\tabcolsep) * \real{0.0937}}
  >{\raggedright\arraybackslash}p{(\columnwidth - 4\tabcolsep) * \real{0.4707}}
  >{\raggedright\arraybackslash}p{(\columnwidth - 4\tabcolsep) * \real{0.4356}}@{}}
\toprule\noalign{}
\begin{minipage}[b]{\linewidth}\raggedright
\end{minipage} & \begin{minipage}[b]{\linewidth}\raggedright
\textbf{Human-like}
\end{minipage} & \begin{minipage}[b]{\linewidth}\raggedright
\textbf{Rational}
\end{minipage} \\
\midrule\noalign{}
\textbf{Thinking} & Cognitive modelling (GPS --- Newell \& Simon) & Laws of thought (logic) \\
\textbf{Acting} & \textbf{Turing test} (1950, "Computing Machinery and Intelligence") & \textbf{Rational agent} (Russell \& Norvig --- the modern approach) \\
\bottomrule\noalign{}
\end{tabular}\hfil\null\par\medskip


\endgroup

\begin{itemize}
\tightlist
\item
  \textbf{Turing test}: interrogator communicates via text with a human and a machine; machine passes if indistinguishable. Needed capabilities: NLP, knowledge representation, automated reasoning, machine learning (+ vision \& robotics for "total Turing test").
\item
  \textbf{Chinese Room} (John Searle): argues passing the test ≠ understanding (strong vs weak AI).
\item
  Term "AI" coined by \textbf{John McCarthy} at \textbf{Dartmouth Conference (1956)}.
\end{itemize}

\needspace{95pt}

\hypertarget{p2-10--11-agents}{%
\subsection{Agents}\label{p2-10--11-agents}}

\diagram{figures/p2-10-00}

\textbf{PEAS} description: Performance measure, Environment, Actuators, Sensors.
(Automated taxi: P = safety, speed, legality; E = roads, traffic; A = steering, brake; S = cameras, GPS, sonar.)

\begingroup\small

\par\medskip\noindent\hfil\begin{tabular}{@{}
  >{\raggedright\arraybackslash}p{(\columnwidth - 2\tabcolsep) * \real{0.1733}}
  >{\raggedright\arraybackslash}p{(\columnwidth - 2\tabcolsep) * \real{0.5822}}@{}}
\toprule\noalign{}
\begin{minipage}[b]{\linewidth}\raggedright
\textbf{Agent type}
\end{minipage} & \begin{minipage}[b]{\linewidth}\raggedright
\textbf{Behaviour}
\end{minipage} \\
\midrule\noalign{}
Simple reflex & condition--action rules on current percept only \\
Model-based reflex & keeps internal state of the world \\
Goal-based & chooses actions to reach goals (search/\allowbreak{}planning) \\
Utility-based & maximises expected utility \\
Learning agent & performance element + learning element + critic + problem generator \\
\bottomrule\noalign{}
\end{tabular}\hfil\null\par\medskip


\endgroup

Environment properties: fully/partially observable, deterministic/stochastic, episodic/sequential, static/dynamic, discrete/continuous, single/multi-agent, known/unknown.
(Chess with clock: fully observable, strategic, sequential, semi-dynamic, discrete, multi-agent.)

\needspace{122pt}

\hypertarget{p2-10--2-state-space-search}{%
\section{State Space Search}\label{p2-10--2-state-space-search}}

Problem = initial state, actions, transition model, \textbf{goal test}, path cost.
Examples: 8-puzzle (9!/2 = 181,440 reachable states), water jug, missionaries \& cannibals, 8-queens, Tower of Hanoi.

\begin{listingblock}{42pt}
\begin{Verbatim}[fontsize=\fontsize{9.0}{10.9}\selectfont]
 Water jug (4 L & 3 L, get 2 L in 4 L jug)
 (0,0) → (0,3) → (3,0) → (3,3) → (4,2) → (0,2) → (2,0)   ✔
\end{Verbatim}
\end{listingblock}

\needspace{160pt}

\hypertarget{p2-10--21-uninformed-search}{%
\subsection{Uninformed Search}\label{p2-10--21-uninformed-search}}

b = branching factor, d = depth of shallowest goal, m = max depth, l = depth limit, C* = optimal cost.

\begingroup\small

\par\medskip\noindent\hfil\begin{tabular}{@{}
  >{\raggedright\arraybackslash}p{(\columnwidth - 8\tabcolsep) * \real{0.3105}}
  >{\raggedright\arraybackslash}p{(\columnwidth - 8\tabcolsep) * \real{0.1957}}
  >{\raggedright\arraybackslash}p{(\columnwidth - 8\tabcolsep) * \real{0.1617}}
  >{\raggedright\arraybackslash}p{(\columnwidth - 8\tabcolsep) * \real{0.2012}}
  >{\raggedright\arraybackslash}p{(\columnwidth - 8\tabcolsep) * \real{0.1265}}@{}}
\toprule\noalign{}
\begin{minipage}[b]{\linewidth}\raggedright
\textbf{Strategy}
\end{minipage} & \begin{minipage}[b]{\linewidth}\raggedright
\textbf{Complete}
\end{minipage} & \begin{minipage}[b]{\linewidth}\raggedright
\textbf{Optimal}
\end{minipage} & \begin{minipage}[b]{\linewidth}\raggedright
\textbf{Time}
\end{minipage} & \begin{minipage}[b]{\linewidth}\raggedright
\textbf{Space}
\end{minipage} \\
\midrule\noalign{}
\textbf{BFS} & Yes (finite b) & Yes (unit cost) & O(b\^{}d) & \textbf{O(b\^{}d)} \\
Uniform-cost & Yes & \textbf{Yes} & O(b\^{}(1+⌊C*/\allowbreak{}ε⌋)) & same \\
\textbf{DFS} & No (infinite paths) & No & O(b\^{}m) & \textbf{O(bm)} \\
Depth-limited & No & No & O(b\^{}l) & O(bl) \\
\textbf{Iterative deepening (IDDFS)} & Yes & Yes (unit cost) & O(b\^{}d) & \textbf{O(bd)} \\
Bidirectional & Yes & Yes & O(b\^{}(d/2)) & O(b\^{}(d/2)) \\
\bottomrule\noalign{}
\end{tabular}\hfil\null\par\medskip


\endgroup

IDDFS combines BFS\textquotesingle s completeness/optimality with DFS\textquotesingle s memory --- preferred uninformed method for large spaces.

\needspace{136pt}

\hypertarget{p2-10--22-heuristic-informed-search}{%
\subsection{Heuristic (Informed) Search}\label{p2-10--22-heuristic-informed-search}}

\begin{listingblock}{86pt}
\begin{Verbatim}[fontsize=\fontsize{9.0}{10.9}\selectfont]
Greedy best-first:  f(n) = h(n)                 — fast, not optimal, not complete
A*:                 f(n) = g(n) + h(n)          — optimal if h is ADMISSIBLE (tree search)
                                                  or CONSISTENT (graph search)
Admissible:   h(n) ≤ h*(n)  (never overestimates)
Consistent:   h(n) ≤ c(n, n') + h(n')    (consistent ⇒ admissible)
Dominance:    h₂ ≥ h₁ (both admissible) ⇒ h₂ better (expands fewer nodes)
\end{Verbatim}
\end{listingblock}

8-puzzle heuristics: h₁ = misplaced tiles, h₂ = \textbf{Manhattan distance} (h₂ dominates h₁).

\textbf{A* worked} --- h(S)=7, h(A)=6, h(B)=2, h(C)=1, h(G)=0 (all admissible)

\diagram{figures/p2-10-01}

\begingroup\small

\par\medskip\noindent\hfil\begin{tabular}{@{}
  >{\raggedright\arraybackslash}p{(\columnwidth - 4\tabcolsep) * \real{0.0636}}
  >{\raggedright\arraybackslash}p{(\columnwidth - 4\tabcolsep) * \real{0.1498}}
  >{\raggedright\arraybackslash}p{(\columnwidth - 4\tabcolsep) * \real{0.6080}}@{}}
\toprule\noalign{}
\begin{minipage}[b]{\linewidth}\raggedright
\textbf{Step}
\end{minipage} & \begin{minipage}[b]{\linewidth}\raggedright
\textbf{Expand (g, f)}
\end{minipage} & \begin{minipage}[b]{\linewidth}\raggedright
\textbf{Open list after expansion (g, f = g + h)}
\end{minipage} \\
\midrule\noalign{}
1 & S (0, 7) & A (1, 7), B (4, 6) \\
2 & B (4, 6) & A (1, 7), C (6, 7) \\
3 & A (1, 7) & \textbf{B (3, 5)} --- cheaper path found, B re-opened; C (6, 7); G (13, 13) \\
4 & B (3, 5) & \textbf{C (5, 6)}; G (13, 13) \\
5 & C (5, 6) & \textbf{G (8, 8)} \\
6 & G (8, 8) & Goal reached: \textbf{S → A → B → C → G, cost 8} (optimal) \\
\bottomrule\noalign{}
\end{tabular}\hfil\null\par\medskip


\endgroup

Note: h here is admissible but \textbf{not consistent} (h(A) = 6 \textgreater{} c(A,B) + h(B) = 4), which is why B had to be re-opened.
With a consistent heuristic, a node is never re-expanded.

\begin{itemize}
\tightlist
\item
  \textbf{AO*} --- for AND-OR graphs (problem decomposition).
\item
  \textbf{IDA*} --- iterative deepening with f-cost limits (memory O(bd)). \textbf{SMA*} --- memory bounded.
\item
  \textbf{Local search}: \textbf{Hill climbing} (steepest ascent) --- problems: \textbf{local maxima, plateaus, ridges}; fixes: random restart, sideways moves. \textbf{Simulated annealing} --- accepts worse moves with probability e\^{}(ΔE/T), T decreases. \textbf{Local beam search}. \textbf{Genetic algorithms} (§8).
\item
  \textbf{Constraint satisfaction (CSP)}: variables, domains, constraints; backtracking + MRV (minimum remaining values), degree heuristic, LCV (least constraining value), forward checking, \textbf{arc consistency (AC-3)}.
\item
  \textbf{Means--ends analysis} (GPS): reduce the difference between current and goal state.
\end{itemize}

\needspace{290pt}

\hypertarget{p2-10--3-game-playing}{%
\section{Game Playing}\label{p2-10--3-game-playing}}

\needspace{254pt}

\hypertarget{p2-10--31-minimax}{%
\subsection{Minimax}\label{p2-10--31-minimax}}

MAX picks maximum, MIN picks minimum child value; complete \& optimal against an optimal opponent. Time O(b\^{}m), space O(bm).

\needspace{188pt}

\hypertarget{p2-10--32-alphabeta-pruning}{%
\subsection{Alpha--Beta Pruning}\label{p2-10--32-alphabeta-pruning}}

α = best value MAX can guarantee so far (lower bound); β = best value MIN can guarantee (upper bound). \textbf{Prune when α ≥ β.}

\begin{listingblock}{108pt}
\begin{Verbatim}[fontsize=\fontsize{9.0}{10.9}\selectfont]
                         MAX                    A = 3
                    ┌─────┴─────┬───────────┐
                   MIN         MIN          MIN
                   B=3         C≤2          D=2
                 ┌─┼─┐       ┌─┼─┐        ┌─┼─┐
                 3 12 8      2  ✂  ✂      14 5 2
                                (pruned: once C sees 2 < α=3, MAX will never choose C)
 Minimax value at root = 3 ; 2 leaves pruned
\end{Verbatim}
\end{listingblock}

\begin{itemize}
\tightlist
\item
  Pruning does \textbf{not} change the final result.
\item
  With perfect move ordering: time \textbf{O(b\^{}(m/2))} (effectively doubles the searchable depth); random ordering ≈ O(b\^{}(3m/4)).
\item
  Stochastic games: \textbf{Expectiminimax} (chance nodes). Imperfect real-time decisions: cutoff test + evaluation function, quiescence search, horizon effect.
\end{itemize}

\needspace{131pt}

\hypertarget{p2-10--4-knowledge-representation}{%
\section{Knowledge Representation}\label{p2-10--4-knowledge-representation}}

\needspace{95pt}

\hypertarget{p2-10--41-logic}{%
\subsection{Logic}\label{p2-10--41-logic}}

\begin{itemize}
\tightlist
\item
  \textbf{Propositional logic}; \textbf{First-order predicate logic (FOPL)} --- see \protect\hyperlink{p2-01--1-mathematical-logic}{Unit 1}.
\item
  Inference: \textbf{resolution} (refutation --- negate goal, convert to \textbf{CNF/clausal form}, derive empty clause), \textbf{forward chaining} (data-driven; production systems, OPS5), \textbf{backward chaining} (goal-driven; Prolog, MYCIN).
\item
  CNF conversion steps: eliminate ↔, → ; move ¬ inward; standardise variables; \textbf{Skolemise} (remove ∃ using Skolem constants/functions); drop ∀; distribute ∨ over ∧.
\item
  \textbf{Unification}: find substitution making literals identical (most general unifier, MGU); occurs check.
\item
  Horn clause: at most one positive literal --- basis of Prolog.
\item
  Monotonic vs \textbf{non-monotonic} reasoning (default logic, closed-world assumption, circumscription, truth-maintenance systems).
\end{itemize}

\needspace{95pt}

\hypertarget{p2-10--42-structured-representations}{%
\subsection{Structured Representations}\label{p2-10--42-structured-representations}}

\diagram{figures/p2-10-02}

\begin{itemize}
\tightlist
\item
  \textbf{Semantic network} (Quillian): nodes = concepts, arcs = relations (is-a, has-a, instance); inheritance.
\item
  \textbf{Frames} (Minsky, 1975): slots \& fillers, default values, procedural attachments (if-needed, if-added demons).
\item
  \textbf{Scripts} (Schank \& Abelson): stereotyped event sequences --- e.g. restaurant script: entry conditions, roles, props, scenes (entering, ordering, eating, exiting), results.
\item
  \textbf{Conceptual Dependency} (Roger Schank): language-independent primitives --- \textbf{ATRANS} (transfer possession: give), \textbf{PTRANS} (physical transfer: go), \textbf{MTRANS} (mental info: tell), \textbf{MBUILD} (build info: decide), \textbf{PROPEL}, \textbf{MOVE}, \textbf{GRASP}, \textbf{INGEST} (eat), \textbf{EXPEL}, \textbf{SPEAK}, \textbf{ATTEND} (11 primitive acts).
\item
  \textbf{Production rules}: IF condition THEN action; conflict resolution (refractoriness, recency, specificity).
\item
  \textbf{Ontologies}: formal specification of a shared conceptualisation (Gruber); OWL, RDF; upper ontologies.
\end{itemize}

\needspace{95pt}

\hypertarget{p2-10--43-expert-systems}{%
\subsection{Expert Systems}\label{p2-10--43-expert-systems}}

\diagram{figures/p2-10-03}

\begingroup\small

\par\medskip\noindent\hfil\begin{tabular}{@{}
  >{\raggedright\arraybackslash}p{(\columnwidth - 2\tabcolsep) * \real{0.2928}}
  >{\raggedright\arraybackslash}p{(\columnwidth - 2\tabcolsep) * \real{0.4822}}@{}}
\toprule\noalign{}
\begin{minipage}[b]{\linewidth}\raggedright
\textbf{Expert system}
\end{minipage} & \begin{minipage}[b]{\linewidth}\raggedright
\textbf{Domain}
\end{minipage} \\
\midrule\noalign{}
\textbf{DENDRAL} (first) & Chemical structure (mass spectrometry) \\
\textbf{MYCIN} & Blood infections --- backward chaining, \textbf{certainty factors} \\
EMYCIN & Shell from MYCIN \\
PROSPECTOR & Mineral exploration \\
XCON / R1 & DEC computer configuration --- forward chaining \\
INTERNIST / CADUCEUS & Internal medicine \\
\bottomrule\noalign{}
\end{tabular}\hfil\null\par\medskip


\endgroup

\needspace{125pt}

\hypertarget{p2-10--44-uncertainty}{%
\subsection{Uncertainty}\label{p2-10--44-uncertainty}}

\begin{listingblock}{75pt}
\begin{Verbatim}[fontsize=\fontsize{9.0}{10.9}\selectfont]
Bayes:   P(H|E) = P(E|H)·P(H) / P(E)
Certainty factor (MYCIN): CF = MB − MD ∈ [−1, 1]
  combining two positive CFs: CF = CF1 + CF2(1 − CF1)
Dempster–Shafer: belief Bel(A) ≤ plausibility Pl(A) ; mass function m over subsets
Bayesian network: DAG + conditional probability tables; P(x₁…xₙ) = Π P(xᵢ | parents(xᵢ))
\end{Verbatim}
\end{listingblock}

Other: fuzzy logic (vagueness --- §7), Markov models, HMMs.

\needspace{130pt}

\hypertarget{p2-10--5-planning}{%
\section{Planning}\label{p2-10--5-planning}}

\begingroup\small

\par\medskip\noindent\hfil\begin{tabular}{@{}
  >{\raggedright\arraybackslash}p{(\columnwidth - 2\tabcolsep) * \real{0.4223}}
  >{\raggedright\arraybackslash}p{(\columnwidth - 2\tabcolsep) * \real{0.5777}}@{}}
\toprule\noalign{}
\begin{minipage}[b]{\linewidth}\raggedright
\textbf{Concept}
\end{minipage} & \begin{minipage}[b]{\linewidth}\raggedright
\textbf{Meaning}
\end{minipage} \\
\midrule\noalign{}
\textbf{STRIPS} (Fikes \& Nilsson, 1971) & Operators with \textbf{preconditions, add list, delete list}; closed-world assumption \\
Linear planning (goal stack) & Solve goals one at a time using a stack --- can fail on interacting goals (\textbf{Sussman anomaly}) \\
Non-linear planning & Goals interleaved; constraint posting \\
\textbf{Partial-order planning (POP)} & Least-commitment; causal links, ordering constraints; resolves \textbf{threats} by promotion/\allowbreak{}demotion \\
\textbf{Hierarchical planning (HTN)} & Abstract actions decomposed into sub-tasks (ABSTRIPS) \\
Graphplan & Planning graph with mutex links \\
Forward (progression) vs backward (regression) state-space planning & \\
\bottomrule\noalign{}
\end{tabular}\hfil\null\par\medskip


\endgroup

\begin{listingblock}{119pt}
\begin{Verbatim}[fontsize=\fontsize{9.0}{10.9}\selectfont]
STRIPS operator (Blocks world)
 PICKUP(x)
   PRE:  ONTABLE(x) ∧ CLEAR(x) ∧ HANDEMPTY
   ADD:  HOLDING(x)
   DEL:  ONTABLE(x), CLEAR(x), HANDEMPTY
 STACK(x, y)
   PRE:  HOLDING(x) ∧ CLEAR(y)
   ADD:  ON(x,y), CLEAR(x), HANDEMPTY
   DEL:  HOLDING(x), CLEAR(y)
\end{Verbatim}
\end{listingblock}

\textbf{Sussman anomaly}: initial C on A, A and B on table; goal ON(A,B) ∧ ON(B,C) --- goal-stack (linear) planning cannot solve it without undoing a subgoal.

\needspace{95pt}

\hypertarget{p2-10--6-natural-language-processing}{%
\section{Natural Language Processing}\label{p2-10--6-natural-language-processing}}

\diagram{figures/p2-10-04}

\begin{itemize}
\tightlist
\item
  Grammars: CFG, \textbf{Transformational grammar} (Chomsky), \textbf{ATN (Augmented Transition Network} --- Woods), Case grammar (Fillmore), DCG, unification grammars, Systemic grammar.
\item
  Parsing: top-down, bottom-up, \textbf{chart parsing} (Earley --- O(n³)), CYK, deterministic parsing.
\item
  Ambiguity: lexical ("bank"), syntactic/structural ("I saw a man with a telescope"), semantic, anaphoric/referential ("he"), pragmatic.
\item
  Semantic analysis: lexical semantics, word sense disambiguation, semantic grammar, case roles.
\item
  \textbf{Pragmatics}: speech acts (Austin, Searle) --- what the speaker intends; \textbf{discourse}: anaphora resolution.
\item
  Modern NLP: n-grams, tokenisation, stemming (Porter) vs lemmatisation, POS tagging (HMM), NER, TF-IDF, word embeddings (word2vec), transformers/LLMs.
\item
  Early systems: \textbf{ELIZA} (Weizenbaum, pattern matching therapist), SHRDLU (Winograd, blocks world), LUNAR, PARRY.
\end{itemize}

\needspace{130pt}

\hypertarget{p2-10--7-multi-agent-systems}{%
\section{Multi-Agent Systems}\label{p2-10--7-multi-agent-systems}}

\begingroup\small

\par\medskip\noindent\hfil\begin{tabular}{@{}
  >{\raggedright\arraybackslash}p{(\columnwidth - 6\tabcolsep) * \real{0.1220}}
  >{\raggedright\arraybackslash}p{(\columnwidth - 6\tabcolsep) * \real{0.2464}}
  >{\raggedright\arraybackslash}p{(\columnwidth - 6\tabcolsep) * \real{0.2798}}
  >{\raggedright\arraybackslash}p{(\columnwidth - 6\tabcolsep) * \real{0.1780}}@{}}
\toprule\noalign{}
\begin{minipage}[b]{\linewidth}\raggedright
\end{minipage} & \begin{minipage}[b]{\linewidth}\raggedright
\textbf{Object}
\end{minipage} & \begin{minipage}[b]{\linewidth}\raggedright
\textbf{Agent}
\end{minipage} & \begin{minipage}[b]{\linewidth}\raggedright
\textbf{Expert system}
\end{minipage} \\
\midrule\noalign{}
Autonomy & Methods invoked by others & Decides itself whether to act & Advises user \\
Behaviour & Passive & Proactive, reactive, social & Passive reasoning \\
Interaction & Method calls & Communication languages & User interface \\
\bottomrule\noalign{}
\end{tabular}\hfil\null\par\medskip


\endgroup

\begin{itemize}
\tightlist
\item
  Agent properties: \textbf{autonomy, reactivity, pro-activeness, social ability} (Wooldridge \& Jennings); BDI architecture (\textbf{Beliefs, Desires, Intentions}).
\item
  MAS structure: agents + environment + interactions (cooperation, coordination, negotiation --- contract net protocol, auctions).
\item
  \textbf{Agent communication languages}: \textbf{KQML} (performatives: ask, tell, achieve\ldots), \textbf{FIPA-ACL}; content language \textbf{KIF}.
\item
  Knowledge sharing via \textbf{ontologies}. \textbf{Semantic web} (Tim Berners-Lee): RDF (triples subject--predicate--object), RDFS, \textbf{OWL}, SPARQL; layered cake.
\item
  Tools: \textbf{JADE} (Java, FIPA-compliant), JACK, Jason (AgentSpeak), Zeus, Aglets, NetLogo.
\end{itemize}

\needspace{166pt}

\hypertarget{p2-10--8-fuzzy-sets--logic-zadeh-1965}{%
\section{Fuzzy Sets \& Logic (Zadeh, 1965)}\label{p2-10--8-fuzzy-sets--logic-zadeh-1965}}

Membership μ\_A(x) ∈ {[}0, 1{]} (vs crisp \{0, 1\}).

\begin{listingblock}{86pt}
\begin{Verbatim}[fontsize=\fontsize{9.0}{10.9}\selectfont]
 μ
 1 ┤      ╱‾‾‾‾╲              triangular(a,b,c): peak at b
   │     ╱      ╲             trapezoidal(a,b,c,d): flat top b..c
   │    ╱        ╲            Gaussian: e^(−(x−c)²/2σ²)
 0 ┼───┴──────────┴───────
       a   b    c   d         Core: μ = 1 ; Support: μ > 0 ; Crossover: μ = 0.5
\end{Verbatim}
\end{listingblock}

\needspace{136pt}

\hypertarget{p2-10--operations-standard-zadeh}{%
\subsection{Operations (standard Zadeh)}\label{p2-10--operations-standard-zadeh}}

\begin{listingblock}{86pt}
\begin{Verbatim}[fontsize=\fontsize{9.0}{10.9}\selectfont]
Union          μ_{A∪B}(x) = max(μA, μB)
Intersection   μ_{A∩B}(x) = min(μA, μB)
Complement     μ_{A'}(x)  = 1 − μA(x)
Algebraic sum  μA + μB − μA·μB ; Algebraic product μA·μB
Bounded sum    min(1, μA + μB) ; Bounded difference max(0, μA − μB)
Concentration (very)   μ²     Dilation (somewhat) √μ
\end{Verbatim}
\end{listingblock}

\textbf{Laws that FAIL in fuzzy sets}: \textbf{law of excluded middle} (A ∪ A\textquotesingle{} ≠ U) and \textbf{law of contradiction} (A ∩ A\textquotesingle{} ≠ ∅). De Morgan, associativity, distributivity still hold.

\textbf{Worked}: A = \{0.2/x₁, 0.7/x₂, 1/x₃\}, B = \{0.5/x₁, 0.3/x₂, 0.8/x₃\}

\begin{listingblock}{53pt}
\begin{Verbatim}[fontsize=\fontsize{9.0}{10.9}\selectfont]
A ∪ B = {0.5, 0.7, 1}      A ∩ B = {0.2, 0.3, 0.8}     A' = {0.8, 0.3, 0}
A ∩ A' = {0.2, 0.3, 0} ≠ ∅
α-cut A₀.₅ = {x₂, x₃}       Height(A) = 1 (normal fuzzy set)
\end{Verbatim}
\end{listingblock}

\begin{itemize}
\tightlist
\item
  \textbf{Fuzzy relations}: max--min composition: (R∘S)(x, z) = max\_y min(R(x,y), S(y,z)). Max-product composition.
\item
  \textbf{Linguistic variables}: e.g. Temperature ∈ \{cold, warm, hot\}; hedges: very, somewhat.
\item
  \textbf{Fuzzy inference system}:
\end{itemize}

\diagram{figures/p2-10-05}

\begin{itemize}
\tightlist
\item
  \textbf{Mamdani} (fuzzy output sets, centroid defuzzification --- intuitive) vs \textbf{Sugeno / TSK} (output = linear function/constant, weighted average --- efficient).
\item
  Centroid: z* = ∫μ(z)·z dz / ∫μ(z) dz. Weighted average: Σ μᵢ zᵢ / Σ μᵢ.
\item
  Applications: washing machines, ABS brakes, air-conditioners, cameras (fuzzy control).
\end{itemize}

\needspace{95pt}

\hypertarget{p2-10--9-genetic-algorithms-john-holland-1975}{%
\section{Genetic Algorithms (John Holland, 1975)}\label{p2-10--9-genetic-algorithms-john-holland-1975}}

\diagram{figures/p2-10-06}

\begin{itemize}
\tightlist
\item
  \textbf{Encoding}: binary, permutation (TSP), value/real, tree (genetic programming).
\item
  \textbf{Fitness function} guides selection. Roulette wheel probability pᵢ = fᵢ / Σf.
\item
  \textbf{Crossover} (exploitation of good genes) with probability \textasciitilde0.6--0.9; \textbf{mutation} (exploration, maintains diversity) \textasciitilde0.001--0.01.
\item
  Single-point crossover: 1011\textbar010 × 0110\textbar111 → 1011111, 0110010.
\item
  \textbf{Schema theorem} (Holland): short, low-order, above-average schemata grow exponentially (building block hypothesis). Order o(H) = fixed positions; defining length δ(H) = distance between first \& last fixed positions. Number of schemata in a string of length l: 3ˡ.
\item
  Premature convergence; elitism keeps the best individuals.
\end{itemize}

\needspace{166pt}

\hypertarget{p2-10--10-artificial-neural-networks}{%
\section{Artificial Neural Networks}\label{p2-10--10-artificial-neural-networks}}

\needspace{130pt}

\hypertarget{p2-10--101-learning-paradigms}{%
\subsection{Learning Paradigms}\label{p2-10--101-learning-paradigms}}

\begingroup\small

\par\medskip\noindent\hfil\begin{tabular}{@{}
  >{\raggedright\arraybackslash}p{(\columnwidth - 4\tabcolsep) * \real{0.1431}}
  >{\raggedright\arraybackslash}p{(\columnwidth - 4\tabcolsep) * \real{0.3483}}
  >{\raggedright\arraybackslash}p{(\columnwidth - 4\tabcolsep) * \real{0.5086}}@{}}
\toprule\noalign{}
\begin{minipage}[b]{\linewidth}\raggedright
\textbf{Type}
\end{minipage} & \begin{minipage}[b]{\linewidth}\raggedright
\textbf{Data}
\end{minipage} & \begin{minipage}[b]{\linewidth}\raggedright
\textbf{Examples}
\end{minipage} \\
\midrule\noalign{}
\textbf{Supervised} & Labelled input--output pairs & Perceptron, back-propagation MLP, SVM, decision trees \\
\textbf{Unsupervised} & Unlabelled & \textbf{SOM (Kohonen)}, k-means, Hebbian, competitive learning, ART \\
\textbf{Reinforcement} & Rewards / penalties from environment & Q-learning, SARSA, TD learning; MDP (states, actions, rewards, policy) \\
\bottomrule\noalign{}
\end{tabular}\hfil\null\par\medskip


\endgroup

Q-learning: Q(s,a) ← Q(s,a) + α{[}r + γ max\_a′ Q(s′,a′) − Q(s,a){]}.

\needspace{114pt}

\hypertarget{p2-10--102-perceptron-rosenblatt-1958}{%
\subsection{Perceptron (Rosenblatt, 1958)}\label{p2-10--102-perceptron-rosenblatt-1958}}

\begin{listingblock}{64pt}
\begin{Verbatim}[fontsize=\fontsize{9.0}{10.9}\selectfont]
 x1 ──w1──╲
 x2 ──w2───► Σ wᵢxᵢ + b ──► step(·) ──► y
 x3 ──w3──╱
 Update rule:  wᵢ ← wᵢ + η (t − y) xᵢ        (t = target, η = learning rate)
\end{Verbatim}
\end{listingblock}

\begin{itemize}
\tightlist
\item
  \textbf{McCulloch--Pitts neuron (1943)}: first model; binary threshold, fixed weights.
\item
  Single-layer perceptron can learn only \textbf{linearly separable} functions --- AND, OR, NAND ✔; \textbf{XOR ✘} (Minsky \& Papert, 1969).
\item
  AND with w₁ = w₂ = 1, threshold θ = 1.5 (or bias −1.5); OR with θ = 0.5.
\item
  Perceptron convergence theorem: converges in finite steps if data is linearly separable.
\end{itemize}

\needspace{95pt}

\hypertarget{p2-10--103-multilayer-perceptron--backpropagation}{%
\subsection{Multilayer Perceptron \& Backpropagation}\label{p2-10--103-multilayer-perceptron--backpropagation}}

\diagram{figures/p2-10-07}

\begin{itemize}
\tightlist
\item
  Input → hidden → output layers; non-linear activations: \textbf{sigmoid} σ(x) = 1/(1 + e⁻ˣ) with σ′ = σ(1 − σ), \textbf{tanh}, \textbf{ReLU} max(0, x), softmax (output).
\item
  \textbf{Backpropagation} (Rumelhart, Hinton, Williams, 1986): forward pass → compute error → propagate gradients backwards (chain rule) → gradient descent: w ← w − η ∂E/∂w. Problems: local minima, vanishing gradients (sigmoid); momentum, learning-rate tuning.
\item
  MLP with one hidden layer is a \textbf{universal approximator}; can solve XOR.
\item
  Overfitting → regularisation, dropout, early stopping.
\item
  \textbf{Delta rule / Widrow--Hoff / LMS} (ADALINE): Δw = η(t − y\_in)x. MADALINE = multiple ADALINEs.
\end{itemize}

\needspace{95pt}

\hypertarget{p2-10--104-self-organising-map-kohonen}{%
\subsection{Self-Organising Map (Kohonen)}\label{p2-10--104-self-organising-map-kohonen}}

\begin{itemize}
\tightlist
\item
  \textbf{Unsupervised, competitive} learning; maps high-dimensional input to a low-dimensional (2-D) grid preserving \textbf{topology}.
\item
  Steps: find \textbf{Best Matching Unit (BMU)} (min Euclidean distance) → update BMU and neighbours: w ← w + η·h(t)(x − w); neighbourhood and η shrink over time.
\end{itemize}

\needspace{95pt}

\hypertarget{p2-10--105-hopfield-network-1982}{%
\subsection{Hopfield Network (1982)}\label{p2-10--105-hopfield-network-1982}}

\begin{itemize}
\tightlist
\item
  \textbf{Recurrent}, fully connected, \textbf{symmetric weights} (wᵢⱼ = wⱼᵢ), \textbf{no self-connections} (wᵢᵢ = 0); binary/bipolar states.
\item
  \textbf{Associative (content-addressable) memory}: stores patterns via Hebbian rule W = Σ pᵀp − I (bipolar); retrieves from noisy input.
\item
  Energy function E = −½ ΣΣ wᵢⱼ sᵢ sⱼ decreases monotonically → converges to stable state (local minimum).
\item
  Capacity ≈ \textbf{0.138 N} patterns for N neurons.
\item
  Also: \textbf{Hebbian learning} ("neurons that fire together wire together") Δw = η·x·y; \textbf{BAM} (bidirectional associative memory --- Kosko); \textbf{ART} (Grossberg -- stability--plasticity); RBF networks; Boltzmann machine.
\end{itemize}

\needspace{183pt}

\hypertarget{p2-10--11-deeper-dive--worked-traces-for-every-numerical-topic}{%
\section{Deeper Dive --- Worked Traces for Every Numerical Topic}\label{p2-10--11-deeper-dive--worked-traces-for-every-numerical-topic}}

\needspace{147pt}

\hypertarget{p2-10--111-alphabeta-trace-depth-3}{%
\subsection{Alpha--Beta Trace (depth 3)}\label{p2-10--111-alphabeta-trace-depth-3}}

\begin{listingblock}{97pt}
\begin{Verbatim}[fontsize=\fontsize{9.0}{10.9}\selectfont]
                         MAX  A
               ┌──────────┴──────────┐
              MIN B                 MIN C
          ┌─────┴─────┐          ┌─────┴─────┐
        MAX D       MAX E      MAX F       MAX G
        ┌─┴─┐       ┌─┴─┐      ┌─┴─┐       ┌─┴─┐
        3   5       6   9      1   2       0  −1
\end{Verbatim}
\end{listingblock}

\begingroup\small

\par\medskip\noindent\hfil\begin{tabular}{@{}
  >{\raggedright\arraybackslash}p{(\columnwidth - 6\tabcolsep) * \real{0.0659}}
  >{\raggedright\arraybackslash}p{(\columnwidth - 6\tabcolsep) * \real{0.0679}}
  >{\raggedright\arraybackslash}p{(\columnwidth - 6\tabcolsep) * \real{0.1559}}
  >{\raggedright\arraybackslash}p{(\columnwidth - 6\tabcolsep) * \real{0.4839}}@{}}
\toprule\noalign{}
\begin{minipage}[b]{\linewidth}\raggedright
\textbf{Step}
\end{minipage} & \begin{minipage}[b]{\linewidth}\raggedright
\textbf{Node}
\end{minipage} & \begin{minipage}[b]{\linewidth}\raggedright
\textbf{α, β on entry}
\end{minipage} & \begin{minipage}[b]{\linewidth}\raggedright
\textbf{What happens}
\end{minipage} \\
\midrule\noalign{}
1 & D & −∞, +∞ & sees 3, 5 → D = 5 \\
2 & B & --- & β = 5 \\
3 & E & −∞, 5 & sees 6 → α = 6 ≥ β = 5 → \textbf{prune leaf 9}; E returns 6 \\
4 & B & --- & B = min(5, 6) = 5 → A\textquotesingle s α = 5 \\
5 & F & 5, +∞ & sees 1, 2 → F = 2 \\
6 & C & 5, +∞ & β = 2 ≤ α = 5 → \textbf{prune G (leaves 0, −1)} \\
7 & A & --- & A = max(5, 2) = \textbf{5} \\
\bottomrule\noalign{}
\end{tabular}\hfil\null\par\medskip


\endgroup

3 of 8 leaves pruned; the root value is the same as plain minimax.

\needspace{155pt}

\hypertarget{p2-10--112-resolution-refutation}{%
\subsection{Resolution Refutation}\label{p2-10--112-resolution-refutation}}

\textbf{First-order}: "All men are mortal. Socrates is a man. Prove Socrates is mortal."

\begin{listingblock}{75pt}
\begin{Verbatim}[fontsize=\fontsize{9.0}{10.9}\selectfont]
1. ¬Man(x) ∨ Mortal(x)          (∀x Man(x) → Mortal(x))
2. Man(Socrates)
3. ¬Mortal(Socrates)            (negated goal)
4. ¬Man(Socrates)               resolve 1, 3 with {x / Socrates}
5. □  (empty clause)            resolve 2, 4  → goal proved
\end{Verbatim}
\end{listingblock}

\textbf{Propositional}: from P → Q, Q → R, P prove R. Clauses ¬P ∨ Q, ¬Q ∨ R, P, ¬R → resolve ¬Q ∨ R with ¬R → ¬Q; with ¬P ∨ Q → ¬P; with P → □.

\needspace{160pt}

\hypertarget{p2-10--113-forward-vs-backward-chaining}{%
\subsection{Forward vs Backward Chaining}\label{p2-10--113-forward-vs-backward-chaining}}

Rules: R1: A ∧ B → C; R2: C → D; R3: D ∧ E → F. Facts: A, B, E. Goal: F.

\begingroup\small

\par\medskip\noindent\hfil\begin{tabular}{@{}
  >{\raggedright\arraybackslash}p{(\columnwidth - 2\tabcolsep) * \real{0.3104}}
  >{\raggedright\arraybackslash}p{(\columnwidth - 2\tabcolsep) * \real{0.3356}}@{}}
\toprule\noalign{}
\begin{minipage}[b]{\linewidth}\raggedright
\textbf{Forward chaining (data-driven)}
\end{minipage} & \begin{minipage}[b]{\linewidth}\raggedright
\textbf{Backward chaining (goal-driven)}
\end{minipage} \\
\midrule\noalign{}
A, B ⊢ C (R1) & Goal F ← need D and E (R3) \\
C ⊢ D (R2) & E is a fact; D ← need C (R2) \\
D, E ⊢ F (R3) ✔ & C ← need A and B (R1); both facts ✔ \\
\bottomrule\noalign{}
\end{tabular}\hfil\null\par\medskip


\endgroup

\needspace{133pt}

\hypertarget{p2-10--114-bayes--medical-test}{%
\subsection{Bayes --- Medical Test}\label{p2-10--114-bayes--medical-test}}

Disease prevalence 1\%, test sensitivity 99\%, false-positive rate 5\%.

\begin{listingblock}{53pt}
\begin{Verbatim}[fontsize=\fontsize{9.0}{10.9}\selectfont]
P(D | +) = P(+ | D)·P(D) / [P(+ | D)·P(D) + P(+ | ¬D)·P(¬D)]
         = 0.99 × 0.01 / (0.0099 + 0.05 × 0.99)
         = 0.0099 / 0.0594 ≈ 0.167      → only about 17% of positives actually have the disease
\end{Verbatim}
\end{listingblock}

\needspace{158pt}

\hypertarget{p2-10--115-fuzzy-maxmin-composition--inference}{%
\subsection{Fuzzy Max--Min Composition \& Inference}\label{p2-10--115-fuzzy-maxmin-composition--inference}}

\begin{listingblock}{108pt}
\begin{Verbatim}[fontsize=\fontsize{9.0}{10.9}\selectfont]
R = | 0.6  0.3 |      S = | 1.0  0.5 |      T = R ∘ S
    | 0.2  0.9 |          | 0.8  0.4 |
T₁₁ = max(min(0.6, 1.0), min(0.3, 0.8)) = max(0.6, 0.3) = 0.6
T₁₂ = max(min(0.6, 0.5), min(0.3, 0.4)) = max(0.5, 0.3) = 0.5
T₂₁ = max(min(0.2, 1.0), min(0.9, 0.8)) = max(0.2, 0.8) = 0.8
T₂₂ = max(min(0.2, 0.5), min(0.9, 0.4)) = max(0.2, 0.4) = 0.4
T = | 0.6  0.5 |
    | 0.8  0.4 |
\end{Verbatim}
\end{listingblock}

\textbf{Sugeno-style inference}: temperature 30 °C → μ\_warm = 0.4, μ\_hot = 0.6. Rules: warm → fan 40\%, hot → fan 80\%.
Output = (0.4 × 40 + 0.6 × 80) / (0.4 + 0.6) = \textbf{64\%}.

\needspace{130pt}

\hypertarget{p2-10--116-genetic-algorithm--one-generation-maximise-fx--xuxb2-5-bit-x}{%
\subsection{Genetic Algorithm --- One Generation (maximise f(x) = x², 5-bit x)}\label{p2-10--116-genetic-algorithm--one-generation-maximise-fx--xuxb2-5-bit-x}}

\begingroup\small

\par\medskip\noindent\hfil\begin{tabular}{@{}
  >{\raggedright\arraybackslash}p{(\columnwidth - 8\tabcolsep) * \real{0.0863}}
  >{\raggedright\arraybackslash}p{(\columnwidth - 8\tabcolsep) * \real{0.0432}}
  >{\raggedright\arraybackslash}p{(\columnwidth - 8\tabcolsep) * \real{0.1074}}
  >{\raggedright\arraybackslash}p{(\columnwidth - 8\tabcolsep) * \real{0.1208}}
  >{\raggedright\arraybackslash}p{(\columnwidth - 8\tabcolsep) * \real{0.2570}}@{}}
\toprule\noalign{}
\begin{minipage}[b]{\linewidth}\raggedright
\textbf{String}
\end{minipage} & \begin{minipage}[b]{\linewidth}\raggedright
\textbf{x}
\end{minipage} & \begin{minipage}[b]{\linewidth}\raggedright
\textbf{f(x)}
\end{minipage} & \begin{minipage}[b]{\linewidth}\raggedright
\textbf{pᵢ = f/Σf}
\end{minipage} & \begin{minipage}[b]{\linewidth}\raggedright
\textbf{Expected copies (4pᵢ)}
\end{minipage} \\
\midrule\noalign{}
01101 & 13 & 169 & 0.144 & 0.58 \\
11000 & 24 & 576 & 0.492 & 1.97 \\
01000 & 8 & 64 & 0.055 & 0.22 \\
10011 & 19 & 361 & 0.309 & 1.23 \\
& & Σ = 1170 & & \\
\bottomrule\noalign{}
\end{tabular}\hfil\null\par\medskip


\endgroup

Crossover 0110\textbar1 × 1100\textbar0 at position 4 → \textbf{01100 (12)} and \textbf{11001 (25)} --- the child 25 is fitter than any parent.

\needspace{130pt}

\hypertarget{p2-10--117-perceptron-learning-and-ux3b7--1-w--0-0-b--0-output-1-if-wx--b--0}{%
\subsection{Perceptron Learning AND (η = 1, w = (0, 0), b = 0, output 1 if w·x + b \textgreater{} 0)}\label{p2-10--117-perceptron-learning-and-ux3b7--1-w--0-0-b--0-output-1-if-wx--b--0}}

\begingroup\small

\par\medskip\noindent\hfil\begin{tabular}{@{}
  >{\raggedright\arraybackslash}p{(\columnwidth - 8\tabcolsep) * \real{0.1499}}
  >{\raggedright\arraybackslash}p{(\columnwidth - 8\tabcolsep) * \real{0.0559}}
  >{\raggedright\arraybackslash}p{(\columnwidth - 8\tabcolsep) * \real{0.0559}}
  >{\raggedright\arraybackslash}p{(\columnwidth - 8\tabcolsep) * \real{0.0469}}
  >{\raggedright\arraybackslash}p{(\columnwidth - 8\tabcolsep) * \real{0.1768}}@{}}
\toprule\noalign{}
\begin{minipage}[b]{\linewidth}\raggedright
\textbf{End of epoch}
\end{minipage} & \begin{minipage}[b]{\linewidth}\raggedright
\textbf{w₁}
\end{minipage} & \begin{minipage}[b]{\linewidth}\raggedright
\textbf{w₂}
\end{minipage} & \begin{minipage}[b]{\linewidth}\raggedright
\textbf{b}
\end{minipage} & \begin{minipage}[b]{\linewidth}\raggedright
\textbf{Errors in epoch}
\end{minipage} \\
\midrule\noalign{}
1 & 1 & 1 & 1 & 1 \\
2 & 2 & 1 & 0 & 3 \\
3 & 2 & 1 & −1 & 3 \\
4 & 2 & 2 & −1 & 2 \\
5 & 2 & 1 & −2 & 1 \\
6 & 2 & 1 & −2 & 0 → converged \\
\bottomrule\noalign{}
\end{tabular}\hfil\null\par\medskip


\endgroup

Check: (0,0) → −2 → 0; (0,1) → −1 → 0; (1,0) → 0 → 0; (1,1) → 1 → 1 ✔ = AND.

\needspace{133pt}

\hypertarget{p2-10--118-one-backpropagation-step-single-sigmoid-neuron}{%
\subsection{One Backpropagation Step (single sigmoid neuron)}\label{p2-10--118-one-backpropagation-step-single-sigmoid-neuron}}

x = 1, w = 0.5, b = 0, target t = 1, η = 1, E = ½(t − o)².

\begin{listingblock}{53pt}
\begin{Verbatim}[fontsize=\fontsize{9.0}{10.9}\selectfont]
net = 0.5 → o = σ(0.5) = 0.6225
δ = (t − o) · o(1 − o) = 0.3775 × 0.6225 × 0.3775 ≈ 0.0887
Δw = η δ x = 0.0887 → w_new = 0.5887 (output moves toward the target)
\end{Verbatim}
\end{listingblock}

\needspace{133pt}

\hypertarget{p2-10--119-hopfield-storage--recall}{%
\subsection{Hopfield Storage \& Recall}\label{p2-10--119-hopfield-storage--recall}}

Store bipolar pattern p = {[}1, −1, 1{]}: W = p pᵀ − I.

\begin{listingblock}{53pt}
\begin{Verbatim}[fontsize=\fontsize{9.0}{10.9}\selectfont]
W = |  0  −1   1 |      Noisy input x = [1, 1, 1]
    | −1   0  −1 |      W x = [0, −2, 0] → sign (keep old value on 0) = [1, −1, 1]
    |  1  −1   0 |      → stored pattern recovered
\end{Verbatim}
\end{listingblock}

\needspace{95pt}

\hypertarget{p2-10--1110-csp--map-colouring-with-heuristics}{%
\subsection{CSP --- Map Colouring with Heuristics}\label{p2-10--1110-csp--map-colouring-with-heuristics}}

\diagram{figures/p2-10-08}

Three colours. \textbf{MRV/degree heuristic} picks SA first (degree 5). Assign SA = red → WA, NT, Q, NSW, V lose red; WA = green → NT = blue → Q = green → NSW = blue → V = green. Solution found without backtracking; \textbf{forward checking} removes inconsistent values early.

\needspace{125pt}

\hypertarget{p2-10--previous-year-questions-pyq-pattern}{%
\section{Previous Year Questions (PYQ pattern)}\label{p2-10--previous-year-questions-pyq-pattern}}

\textbf{Foundations \& agents}

\begin{enumerate}
\def\labelenumi{\arabic{enumi}.}
\tightlist
\item
  The term "Artificial Intelligence" was coined by: \textbf{John McCarthy (1956)}
\item
  Turing test evaluates: \textbf{whether a machine\textquotesingle s behaviour is indistinguishable from a human\textquotesingle s}
\item
  PEAS stands for: \textbf{Performance, Environment, Actuators, Sensors}
\item
  Agent that acts only on the current percept: \textbf{simple reflex agent}
\item
  Chess is: \textbf{fully observable, deterministic (strategic), sequential, discrete, multi-agent}
\end{enumerate}

\textbf{Search}

\begin{enumerate}
\def\labelenumi{\arabic{enumi}.}
\setcounter{enumi}{5}
\tightlist
\item
  Space complexity of BFS: \textbf{O(b\^{}d)}; of DFS: \textbf{O(bm)}
\item
  Search that combines benefits of BFS and DFS: \textbf{iterative deepening DFS}
\item
  A* is optimal if the heuristic is: \textbf{admissible (never overestimates)}
\item
  If h(n) = 0 for all n, A* reduces to: \textbf{uniform-cost search} (Dijkstra)
\item
  Greedy best-first search uses: \textbf{f(n) = h(n)}
\item
  Hill climbing can get stuck due to: \textbf{local maxima, plateau, ridges}
\item
  Simulated annealing accepts worse moves with probability: \textbf{e\^{}(ΔE/T)}
\item
  Problem reduction (AND-OR graphs) uses: \textbf{AO*}
\item
  Means-ends analysis was used in: \textbf{GPS (General Problem Solver)}
\item
  For 8-puzzle, Manhattan distance is: \textbf{admissible and dominates misplaced tiles}
\end{enumerate}

\textbf{Games}

\begin{enumerate}
\def\labelenumi{\arabic{enumi}.}
\setcounter{enumi}{15}
\tightlist
\item
  Alpha-beta pruning with perfect ordering reduces time to: \textbf{O(b\^{}(m/2))}
\item
  Pruning occurs when: \textbf{α ≥ β}
\item
  Alpha-beta pruning changes the minimax value? \textbf{No}
\item
  Minimax value of tree with MIN nodes having leaves (3,12,8), (2,4,6), (14,5,2): max(3, 2, 2) = \textbf{3}
\end{enumerate}

\textbf{Knowledge representation}

\begin{enumerate}
\def\labelenumi{\arabic{enumi}.}
\setcounter{enumi}{19}
\tightlist
\item
  Frames were proposed by: \textbf{Marvin Minsky}
\item
  Scripts were proposed by: \textbf{Schank \& Abelson}
\item
  "John gave Mary a book" in CD uses primitive: \textbf{ATRANS}
\item
  MYCIN used: \textbf{backward chaining with certainty factors}
\item
  First expert system: \textbf{DENDRAL}
\item
  Skolemisation removes: \textbf{existential quantifiers}
\item
  Prolog is based on: \textbf{Horn clauses with backward chaining (SLD resolution)}
\item
  Resolution proves a goal by: \textbf{refutation (deriving the empty clause from negated goal)}
\item
  Combining CF 0.6 and 0.4 (both positive): 0.6 + 0.4 × 0.4 = \textbf{0.76}
\item
  Dempster--Shafer theory uses: \textbf{belief and plausibility}
\end{enumerate}

\textbf{Planning \& NLP}

\begin{enumerate}
\def\labelenumi{\arabic{enumi}.}
\setcounter{enumi}{29}
\tightlist
\item
  STRIPS operators contain: \textbf{preconditions, add list, delete list}
\item
  Sussman anomaly illustrates a problem with: \textbf{linear (goal-stack) planning}
\item
  Partial-order planning follows: \textbf{least-commitment strategy}
\item
  ATN stands for: \textbf{Augmented Transition Network}
\item
  "I saw the man with the telescope" exhibits: \textbf{structural (syntactic) ambiguity}
\item
  Correct NLP pipeline: \textbf{morphological → syntactic → semantic → discourse → pragmatic}
\item
  ELIZA was developed by: \textbf{Joseph Weizenbaum}
\end{enumerate}

\textbf{MAS}

\begin{enumerate}
\def\labelenumi{\arabic{enumi}.}
\setcounter{enumi}{36}
\tightlist
\item
  BDI stands for: \textbf{Beliefs, Desires, Intentions}
\item
  KQML is used for: \textbf{agent communication}
\item
  JADE is: \textbf{a Java framework for FIPA-compliant multi-agent systems}
\item
  Semantic web ontology language: \textbf{OWL}
\end{enumerate}

\textbf{Fuzzy}

\begin{enumerate}
\def\labelenumi{\arabic{enumi}.}
\setcounter{enumi}{40}
\tightlist
\item
  Fuzzy logic was introduced by: \textbf{Lotfi Zadeh (1965)}
\item
  μA = 0.6, μB = 0.3: union \textbf{0.6}, intersection \textbf{0.3}, complement of A \textbf{0.4}
\item
  Which law does NOT hold in fuzzy sets? \textbf{Law of excluded middle / contradiction}
\item
  Defuzzification method computing centre of area: \textbf{centroid}
\item
  Sugeno fuzzy model output is: \textbf{a linear function or constant (crisp)}
\item
  Max--min composition of R = {[}0.3 0.8{]} and S = {[}0.5; 0.9{]} (column): max(min(0.3,0.5), min(0.8,0.9)) = \textbf{0.8}
\item
  "Very" hedge applied to μ = 0.8: \textbf{0.64}
\end{enumerate}

\textbf{GA}

\begin{enumerate}
\def\labelenumi{\arabic{enumi}.}
\setcounter{enumi}{47}
\tightlist
\item
  Genetic algorithms were introduced by: \textbf{John Holland}
\item
  Operator that maintains diversity: \textbf{mutation}
\item
  Roulette wheel: fitness values 10, 20, 30, 40 → probability of 3rd: \textbf{0.3}
\item
  Number of schemata in binary string length 5: \textbf{3⁵ = 243}
\end{enumerate}

\textbf{Neural networks}

\begin{enumerate}
\def\labelenumi{\arabic{enumi}.}
\setcounter{enumi}{51}
\tightlist
\item
  Single-layer perceptron cannot learn: \textbf{XOR}
\item
  Perceptron learning rule: \textbf{Δw = η(t − y)x}
\item
  Backpropagation uses: \textbf{gradient descent with chain rule}
\item
  Derivative of sigmoid: \textbf{σ(1 − σ)}
\item
  Kohonen SOM is: \textbf{unsupervised, competitive learning preserving topology}
\item
  Hopfield network weights are: \textbf{symmetric with zero diagonal}; used as \textbf{associative memory}
\item
  Hopfield storage capacity: \textbf{≈ 0.138 N}
\item
  Learning with reward/penalty: \textbf{reinforcement learning}
\item
  Hebbian learning rule: \textbf{Δw = η·x·y}
\end{enumerate}

\textbf{More practice questions}

\begin{enumerate}
\def\labelenumi{\arabic{enumi}.}
\setcounter{enumi}{60}
\tightlist
\item
  In the tree of §11.1, the number of leaves pruned by alpha--beta: \textbf{3}
\item
  Resolution proves a goal by deriving: \textbf{the empty clause from the negated goal}
\item
  Backward chaining starts from: \textbf{the goal}
\item
  Prevalence 1\%, sensitivity 99\%, false-positive 5\%: P(disease \textbar{} positive) ≈ \textbf{0.17}
\item
  Max--min composition entry max(min(0.2, 1.0), min(0.9, 0.8)) = \textbf{0.8}
\item
  In a GA with fitness values 169, 576, 64, 361, the selection probability of the fittest: \textbf{≈ 0.49}
\item
  Single-point crossover of 01101 and 11000 after bit 4: \textbf{01100 and 11001}
\item
  Perceptron weights (w₁, w₂, b) = (2, 1, −2) implement: \textbf{AND}
\item
  Derivative of the sigmoid at output 0.5: \textbf{0.25}
\item
  Hopfield weight matrix for one stored pattern p: \textbf{p pᵀ − I}
\item
  In map colouring, choosing the variable with the fewest legal values is the: \textbf{MRV heuristic}
\item
  Degree heuristic chooses the variable: \textbf{involved in the most constraints with unassigned variables}
\end{enumerate}

\begin{revbox}

\begin{itemize}
\tightlist
\item
  BFS O(b\^{}d) space · DFS O(bm) · IDDFS best uninformed · A* optimal if admissible
\item
  α--β: prune when α ≥ β · best ordering O(b\^{}(m/2))
\item
  Minsky frames · Schank scripts \& CD (ATRANS, PTRANS, MTRANS) · MYCIN backward + CF
\item
  STRIPS: pre/add/delete · Sussman anomaly → linear planning fails
\item
  Fuzzy: ∪ max, ∩ min, \textquotesingle{} = 1 − μ; excluded middle fails
\item
  Perceptron no XOR · SOM unsupervised · Hopfield symmetric, recurrent, 0.138N
\end{itemize}

\end{revbox}
